\begin{table}[htbp]
\centering
\caption{Matched RGB-occlusion comparison on Unity, with supplied calibration and no added extrinsic perturbation. Common-13 canonical MPJPE (lower is better), on all 64,440 windows. Ratios select joints for patch placement, not image-area coverage.}
\label{tab:c07_occlusion_mpjpe}
\small
\setlength{\tabcolsep}{4pt}
\begin{tabular}{lrrrrr}
\toprule
Condition & CanonFuse3D & Average & DLT & Gated DLT & DLT-residual \\
\midrule
Clean & 0.1553 & 0.2797 & 0.2004 & 0.2108 & 0.0988 \\
Left 0.5 & 0.2103 & 0.3051 & 0.2542 & 0.2606 & 0.1498 \\
Right 0.5 & 0.1949 & 0.2890 & 0.2488 & 0.2556 & 0.1417 \\
Both 0.5 & 0.2514 & 0.3146 & 0.2957 & 0.2983 & 0.1870 \\
Left 1.0 & 0.4654 & 0.3992 & 0.4499 & 0.4495 & 0.3569 \\
Right 1.0 & 0.3706 & 0.3347 & 0.5501 & 0.4569 & 0.4070 \\
Both 1.0 & 0.5986 & 0.4654 & 0.5121 & 0.4866 & 0.4155 \\
\bottomrule
\end{tabular}
\par\smallskip\begin{minipage}{\linewidth}\footnotesize CanonFuse3D and canonical averaging use only predicted 3D. The three DLT-based methods additionally use estimated 2D and calibration. DLT follows~\cite{hartley2004multiple}; gating and the residual MLP are study-specific controls, not reproductions of published learned triangulation systems.\end{minipage}
\end{table}

